% Generated draft: requires mathematical and source review.
\documentclass[11pt]{article}
\usepackage[margin=25mm]{geometry}
\usepackage{fontspec}
\setmainfont{Latin Modern Roman}
\usepackage{amsmath,amssymb,mathtools,unicode-math}
\setmathfont{Latin Modern Math}
\usepackage{enumitem,needspace,xurl,hyperref}
\setlength{\parindent}{0pt}
\setlength{\parskip}{5pt}
\setlength{\emergencystretch}{3em}
\newcommand{\N}{\mathbb{N}}
\newcommand{\Z}{\mathbb{Z}}
\newcommand{\Q}{\mathbb{Q}}
\newcommand{\R}{\mathbb{R}}
\newcommand{\C}{\mathbb{C}}
\newcommand{\F}{\mathbb{F}}
\newcommand{\A}{\mathbb{A}}
\newcommand{\PP}{\mathbb{P}}
\newcommand{\E}{\mathbb{E}}
\newcommand{\eps}{\varepsilon}
\newcommand{\dd}{\,\mathrm{d}}
\begin{document}
\section*{UnsolvedMath ID: 30002618}
\textbf{Source problem code:} catalogue entry
\subsection{Definitions and mathematical statement}
The source considers a semistable curve, denoted through its special fibre by \(X_k\), and an \(F\)-convergent isocrystal \(E\). Here an \(F\)-convergent isocrystal means a convergent isocrystal equipped with a Frobenius structure. The source does not specify the base, the residue field \(k\), the coefficient field, properness assumptions, or the logarithmic structures. These data must be supplied before the cohomology groups and comparison maps below can be defined precisely.

Write
\[
V_E=H^1_{\mathrm{rig}}(X_k,E),\qquad
W_E=H^1_{\mathrm{log\text{-}cris}}(X_k,E),
\]
for the degree-one rigid and logarithmic crystalline cohomology groups with coefficients indicated by \(E\). Let \(\alpha_E:V_E\to W_E\) denote the unnamed comparison map intended in the source, and let \(N_E\) denote its intended monodromy operator. The notation presupposes compatible coefficient objects for both cohomology theories; the source does not specify how these are obtained from \(E\).

\textbf{Problem (conditional on specifying these data).} Determine which coefficient objects \(E\) admit the intended maps and satisfy exactness at the middle term of
\[
H^1_{\mathrm{rig}}(X_k,E)
\xrightarrow{\alpha_E}
H^1_{\mathrm{log\text{-}cris}}(X_k,E)
\xrightarrow{N_E}
H^1_{\mathrm{log\text{-}cris}}(X_k,E).
\]
Explicitly, classify those \(E\) for which
\[
\operatorname{im}(\alpha_E)=\ker(N_E).
\]
Thus every element killed by \(N_E\) must come from rigid cohomology, and every element coming from rigid cohomology must be killed by \(N_E\). No injectivity of \(\alpha_E\) or surjectivity of \(N_E\) is requested by this displayed sequence.

The target of \(N_E\) is reproduced exactly as supplied; the source does not explain any Frobenius twist or whether Frobenius structures are being forgotten. It also does not specify whether the desired classification is for each fixed curve or uniformly over a class of curves. The supplied background particularly emphasizes coefficients arising geometrically, without defining that class. Consequently, the equality above records the intended exactness criterion, while the ambient classification problem remains dependent on unresolved hypotheses.
\subsection{Short English statement}
Determine which Frobenius-equipped convergent isocrystals on a semistable curve make the image of rigid cohomology equal to the kernel of monodromy on logarithmic crystalline cohomology, with particular interest in geometrically arising coefficients. The base data, compatible coefficient constructions, maps, and scope of the classification require clarification.
\subsection{Sources}
\begin{itemize}
\item Local UnsolvedMath registry, permanent ID \(30002618\), source code \texttt{catalogue entry}, supplied complete source record; registry locations: \texttt{data/unsolvedmath/problems.json} and \texttt{catalog/problems.json}. Title: \emph{Monodromy Exactness for Isocrystals on Semistable Curves}. No individual proposer is supplied.
\item Supplied report attribution: \emph{Algebraische Zahlentheorie (2015)}, Oberwolfach Reports, pp.~1747--1800, DOI \texttt{10.4171/owr/2014/32}. \url{https://doi.org/10.4171/owr/2014/32}. This is also the supplied literature source labelled ``original Oberwolfach report.'' The registry gives proposed year 2014; the supplied report title gives 2015, while the DOI contains 2014. This dating discrepancy is unresolved.
\item Supplied literature source: Chiarellotto--Coleman--Di Proietto--Iovita, \emph{p-adic invariant cycles}. \url{https://doi.org/10.1515/crelle-2013-0117}. No theorem number or precise supporting passage is supplied.
\item Literature assessment dated \texttt{2026-08-21}: the supplied curation labels the problem open and states that non-exactness is known for general coefficients, while classification of coefficients giving exactness, especially geometrically arising coefficients, remains open. This is supplied dated curation, not an independently verified literature assessment; its status and cited claims require independent verification.
\end{itemize}
\end{document}
